% !TEX program = pdflatex
% Documentation fonction par fonction de graph_pipeline_prototype.ipynb
% Compilation : pdflatex graph_pipeline_documentation.tex (2 passes pour les références)

\documentclass[11pt,a4paper]{article}

\usepackage[utf8]{inputenc}
\usepackage[T1]{fontenc}
\usepackage[french]{babel}
\usepackage{lmodern}
\usepackage[margin=2.2cm]{geometry}
\usepackage{booktabs}
\usepackage{array}
\usepackage{longtable}
\usepackage{xcolor}
\usepackage[hidelinks]{hyperref}

% Désactive les espacements automatiques du français devant ; : ! ?
\shorthandoff*{;:!?}

\newcommand{\code}[1]{\texttt{\small #1}}

\title{Documentation de \code{graph\_pipeline\_prototype.ipynb}}
\author{}
\date{}

\begin{document}
\maketitle

%======================================================================
\section{Hypothèses métier}
\label{sec:hypotheses}

{\small
\begin{longtable}{@{}l p{2.45cm} p{5.9cm} p{4.5cm}@{}}
\toprule
\textbf{Id} & \textbf{Étape} & \textbf{Hypothèse métier} & \textbf{Si elle est fausse} \\
\midrule
\endfirsthead
\toprule
\textbf{Id} & \textbf{Étape} & \textbf{Hypothèse métier} & \textbf{Si elle est fausse} \\
\midrule
\endhead
\bottomrule
\endfoot

H1 & \url{get_customer} & Le \code{shortName} désigne un client unique, qui possède un publisher dans le namespace \code{artemis-prod}. & La jointure \code{inner} rend le client invisible, ou \code{collect()} renvoie plusieurs \code{db\_name} et le graphe mélange deux catalogues. \\
\addlinespace
H2 & \url{get_products} & Le catalogue \emph{courant} vaut pour toute la fenêtre historique : un produit actif aujourd'hui l'était déjà lors des vues de janvier. & Les produits désactivés depuis sont retirés du graphe alors que leurs vues existaient : le graphe rétrécit à mesure que l'on relance tard. \\
\addlinespace
H3 & \url{get_products} & Seuls les produits maîtres (\code{parentId = -1}) portent le signal ; le tracking a déjà ramené les vues de déclinaisons au maître. & Le signal des déclinaisons est perdu, ou au contraire agrégé alors que les déclinaisons se comportent différemment. \\
\addlinespace
H4 & \url{get_products} & \code{isActive} $\land$ \code{isOkInStream} = « produit recommandable » ; les autres n'ont pas à être des nœuds. & Des produits recommandés en production sont absents du graphe, donc non scorables. \\
\addlinespace
H5 & \url{get_embeddings} & Un modèle d'embedding figé (\code{last\_model\_id}) donne au plus une ligne par produit. & Deux embeddings pour le même produit se retrouvent tous les deux dans le résultat, et le remplissage des features en aval choisit l'un des deux de façon non déterministe. \\
\addlinespace
H6 & \url{get_views} & Un événement \code{product\_display\_page\_event} = une vue produit. Pas de pondération par durée ou scroll, pas de déduplication des rafraîchissements. & Un rechargement de page compte comme une vue supplémentaire ; les vues de bots gonflent les sessions. \\
\addlinespace
H7 & \url{get_views} & \code{t2s\_merged\_user} est fonctionnel (un \code{sourceId} $\rightarrow$ un \code{masterId}) et l'identité réconciliée est la bonne granularité de session. & Plusieurs \code{masterId} pour un \code{sourceId} dupliquent les lignes de vues et \emph{doublent} les poids d'arêtes. Fusionner les appareils enchaîne des navigations parallèles dans une même session. \\
\addlinespace
H8 & \url{get_views} & Trafic non identifié (\code{user.internalId} null) et vue non résolue en produit maître ne portent pas de signal exploitable. & Une part importante du trafic --- typiquement les visiteurs non logués --- ne contribue à aucune arête. \\
\addlinespace
H9 & \url{get_views} & \code{emitted.between(start\_date, end\_date)} : les bornes sont des dates converties en timestamps à minuit, donc la journée du \code{end\_date} est de fait exclue. & La fenêtre réelle a un jour de moins que prévu ; un décalage d'un jour entre le graphe et un autre calcul qui utiliserait la même variable. \\
\addlinespace
H10 & \url{sessionize_views} & Une session = 30 minutes d'inactivité (\code{session\_timeout}), par utilisateur réconcilié, tous appareils confondus. Le tri \code{(emitted, internalId)} rend la coupure déterministe en cas d'\emph{ex æquo}. & Un temps de réflexion long scinde une intention d'achat en deux sessions ; deux intentions rapprochées fusionnent. \\
\addlinespace
H11 & \url{split_sessions} & Le découpage train/val/test se fait à la granularité de la session (\code{session\_start}), pas de l'événement individuel. & Une session commencée avant \code{cutoff\_train} et poursuivie après reste entièrement en train : de l'information postérieure à la coupure entre dans l'entraînement. \\
\addlinespace
H12 & \url{build_co_view_edges} & Deux produits vus dans la même session sont liés, indépendamment de l'ordre, de l'écart de temps et de la distance dans la session. & Une session de navigation exploratoire produit un graphe complet de produits sans rapport ; le signal séquentiel (A puis B) est perdu. \\
\addlinespace
H13 & \url{build_co_view_edges} & Une session compte une fois par paire (\code{distinct} avant la self-jointure) : le poids d'une arête = nombre de sessions distinctes contenant la paire. & Un utilisateur qui compare intensément deux produits ne pèse pas plus qu'une co-vue fortuite au sein d'une session. \\
\addlinespace
H14 & \url{build_co_view_edges} & \code{USE\_SESSION\_WEIGHT\_NORM = False} : toutes les sessions pèsent pareil. Une session de $n$ produits produit $n(n-1)/2$ paires de poids 1. & Une session de 50 produits (bot, crawler, comparateur) apporte 1225 arêtes contre 1 pour une session de 2 produits : le graphe est dominé par les sessions longues. \\
\addlinespace
H15 & \url{build_co_view_edges} & Aucun seuil minimal de poids : une co-occurrence unique est une arête à part entière. & Le graphe est majoritairement constitué d'arêtes de poids 1, c'est-à-dire de coïncidences --- et ce sont exactement celles qui varient le plus d'un run à l'autre. \\
\addlinespace
H16 & \url{build_co_view_edges} & \code{L.internalId < R.internalId} élimine les boucles et canonise chaque paire dans un seul sens. & Aucune asymétrie de type « B est un accessoire de A » n'est représentable ici ; la duplication dans les deux sens n'a lieu qu'en aval. \\
\addlinespace
H17 & \url{remove_train_edges} & Une paire évaluée doit être inédite vis-à-vis du graphe d'entraînement. & Sans cela on mesurerait la mémorisation plutôt que la généralisation. Val et test ne sont pas dédoublonnés entre eux : une même paire peut être dans les deux. \\
\addlinespace
H18 & \url{build} & \code{df\_train\_edges is None} sert de marqueur implicite « ceci est le split train » ; c'est le seul signal qui déclenche ou non le retrait des paires déjà en train. & Un appel val/test sans passer \code{df\_train\_edges} ne filtre silencieusement rien : aucune vérification n'assure que le split réellement traité correspond à l'usage prévu du paramètre. \\
\addlinespace
H19 & \url{build_node_mapping} & Le jeu de nœuds est l'union des arêtes de train, val et test passées en argument. & Des nœuds existent dans le graphe sans aucune arête d'entraînement (isolés), uniquement parce qu'ils ont été co-vus pendant val ou test. Ils augmentent le nombre de nœuds et donc la taille des couches d'embedding. \\
\addlinespace
H20 & \url{build_node_mapping} & Les index de nœuds sont attribués dans l'ordre croissant des \code{internalId}. & Reproductible à jeu de nœuds identique ; mais l'ajout d'un seul produit décale tous les index suivants et rend un fichier de données incompatible avec un \code{node2idx.json} d'un autre run. \\
\addlinespace
H21 & \url{remap_edges} & \code{node2idx} couvre déjà tous les \code{internalId} de \code{df}, car construit en amont sur ce même DataFrame (entre autres). & \code{.map()} renvoie \code{NaN} pour un \code{internalId} absent, silencieusement --- la ligne n'est pas retirée mais devient corrompue, et l'erreur ne surviendrait qu'en aval, à la conversion en tensor. \\
\addlinespace
H22 & \url{edges_to_tensors} & Le graphe de co-vue est traité comme non dirigé : \code{symmetric=True} ajoute systématiquement l'arête inverse, complétant la canonisation \code{A < B} faite dans \code{build\_co\_view\_edges}. & Un appel avec \code{symmetric=False} produirait un graphe dirigé avec deux fois moins d'arêtes, sans aucun avertissement puisque le paramètre n'est validé contre rien. \\
\addlinespace
H23 & \url{recommended_products_for_t2s_user} & Un log ad et une vue produit du même cookie, du même produit et du même client, séparés de moins de 3 secondes, sont le même événement. & Fenêtre trop large : plusieurs vues apparient un même log et dupliquent les lignes. Trop étroite : un log est perdu par latence réseau ou décalage d'horloge. \\
\addlinespace
H24 & \url{recommended_products_for_t2s_user} & Seules les pages \code{PRODUCT} portent une recommandation exploitable. & Les recommandations de la page d'accueil, de la recherche et des pages catégorie sont hors périmètre. \\
\addlinespace
H25 & \url{positive_negative_edge_construction} & L'ordre du tableau \code{relevantProducts} est l'ordre d'affichage réel, et seuls les 12 premiers comptent (\code{F.slice(...,1,12)}). & Le code prévoit explicitement une variante par score (\code{relevanceScore}) ailleurs dans le notebook ; trier autrement changerait quels candidats entrent dans les 12 retenus. \\
\addlinespace
H26 & \url{positive_negative_edge_construction} & \textbf{Positif} = candidat recommandé \emph{et} vu dans la même session --- sans contrainte temporelle : un produit vu avant la recommandation compte comme positif. & Fuite temporelle dans la définition du label : une partie des « clics prédits » sont des produits déjà visités, que le modèle peut retrouver par le graphe plutôt que par une vraie prédiction. \\
\addlinespace
H27 & \url{positive_negative_edge_construction} & \textbf{Négatif} = candidat recommandé, non positif, \emph{et} absent des arêtes d'entraînement testées dans les deux sens. & Évite d'apprendre contre des faux négatifs, mais retire les négatifs les plus difficiles --- ceux qui ressemblent le plus à des positifs. \\
\addlinespace
H28 & \url{positive_negative_edge_construction} & Asymétrie assumée : les négatifs sont filtrés sur les arêtes du graphe d'entraînement, les positifs ne le sont pas. & Un modèle qui apprend seulement « cette paire est-elle déjà une arête du graphe ? » obtient un bon score sans généraliser. \\
\addlinespace
H29 & \url{positive_negative_edge_construction} & Négatifs tronqués au rang du positif le plus profond de l'exécution ; les exécutions sans aucun positif sont écartées des arêtes d'entraînement (\code{sample(fraction=0.0, seed=42)}). & Le modèle ne voit jamais de liste purement négative, alors que la production en produit. La fraction à 0 est un levier laissé tel quel, pas une propriété du problème. \\
\addlinespace
H30 & \url{positive_negative_edge_construction} & \code{exec2code} couvre toutes les exécutions du split, indépendamment du label, et est rendu stable par tri (\code{all\_exec\_ids.sort()}). & Stable à jeu d'exécutions identique seulement : la moindre exécution en plus ou en moins renumérote tout et invalide les tenseurs déjà écrits avec l'ancien \code{exec2code}. \\
\addlinespace
H31 & \url{prod_products_by_trigger} & Aujourd'hui, chaque trigger n'a qu'un seul placement sponsorisé ; le \code{rank} final coïncide donc exactement avec la position dans le tableau (1-based). & Si un trigger avait plusieurs placements, la fusion par \code{F.min(pos\_rank)} et le tie-break par \code{internal\_id} entreraient en jeu silencieusement, sans qu'aucun test ne le signale aujourd'hui. \\
\addlinespace
H32 & \url{add_pos_neg} & Pour la comparaison à la production, positif et négatif sont définis par simple appartenance aux produits de la session, sans aucun filtrage sur les arêtes du graphe d'entraînement. & Choix assumé : on veut le vrai retour production, pas des négatifs d'entraînement. Conséquence : la baseline production et le modèle ne sont pas évalués sur exactement la même définition de négatif que celle utilisée à l'entraînement (H26). \\

\end{longtable}
}
\refstepcounter{table}\label{tab:hypotheses}
\begin{center}\footnotesize\itshape Tableau \thetable{} --- Hypothèses métier.\end{center}

%======================================================================
\clearpage
\section{Ce qui peut être exclu du graphe}
\label{sec:exclusions}

{\small
\begin{longtable}{@{}l p{3.2cm} p{4.4cm} p{5.2cm}@{}}
\toprule
\textbf{Id} & \textbf{Étape} & \textbf{Mécanisme} & \textbf{Ce qui disparaît} \\
\midrule
\endfirsthead
\toprule
\textbf{Id} & \textbf{Étape} & \textbf{Mécanisme} & \textbf{Ce qui disparaît} \\
\midrule
\endhead
\bottomrule
\endfoot

E1 & \url{get_customer} & jointure \code{inner} clients / publishers. & Un client sans \emph{publisher} \code{artemis-prod} : le notebook s'exécute alors sur des listes vides sans erreur. \\
\addlinespace
E2 & \url{get_products} & \code{isActive}, \code{isOkInStream} & Produits désactivés ou hors flux au moment du run (pas au moment de la vue). \\
\addlinespace
E3 & \url{get_products} & \code{parentId = -1} & Toutes les déclinaisons. \\
\addlinespace
E4 & \url{get_embeddings} & \code{last\_model\_id} figé & Les embeddings calculés par une autre version du modèle. \\
\addlinespace
E5 & \url{get_views} & \code{user.internalId} non nul & Trafic non identifié. \\
\addlinespace
E6 & \url{get_views} & \code{masterProductInternalId} non nul & Vues non résolues en produit maître. \\
\addlinespace
E7 & \url{get_views} & \code{emitted.between(...)} & Journée du \code{end\_date} (bornes prises à minuit, voir H9). \\
\addlinespace
E8 & \url{build_co_view_edges} & self-jointure \code{L.internalId < R.internalId} & Sessions à un seul produit distinct : elles ne produisent aucune arête, quel que soit leur volume. \\
\addlinespace
E9 & \url{remove_train_edges} & \code{left\_anti} & Paires val/test déjà vues en train. \\
\addlinespace
E10 & \url{build_node_mapping} & \code{node2idx} n'est construit qu'à partir des colonnes \code{product\_A}/\code{product\_B} des arêtes fournies. & Les produits catalogués et embeddés qui n'ont jamais été co-vus avec un autre produit ne deviennent jamais des nœuds du graphe, même s'ils ont un embedding valide. \\
\addlinespace
E11 & \url{recommended_products_for_t2s_user} & \code{pageType = "PRODUCT"}, fenêtre $\pm$3\,s, \code{publisherId isin(...)} & Recommandations hors page produit ; logs non appariés à une vue ; logs d'un publisher non résolu. \\
\addlinespace
E12 & \url{positive_negative_edge_construction} & \code{F.explode(...)} sur le tableau des placements sponsorisés & Exécutions sans placement sponsorisé (tableau vide ou nul) : \code{explode} les retire silencieusement, contrairement à \code{explode\_outer}. \\
\addlinespace
E13 & \url{positive_negative_edge_construction} & \code{F.slice(relevantProducts, 1, 12)} & Candidats au-delà du rang 12 --- y compris d'éventuels positifs qui s'y trouveraient. \\
\addlinespace
E14 & \url{positive_negative_edge_construction} & \code{\textasciitilde{}is\_existing\_edge} & Paires candidates déjà arêtes du graphe d'entraînement, côté négatifs uniquement (H28). \\
\addlinespace
E15 & \url{positive_negative_edge_construction} & \code{rank\_candidate <= deepest\_pos\_rank} & Négatifs plus profonds que le positif le plus profond de l'exécution. \\
\addlinespace
E16 & \url{positive_negative_edge_construction} & \code{sample(fraction=0.0, seed=42)} & Tous les négatifs des exécutions sans aucun positif. \\
\addlinespace
E17 & \url{positive_negative_edge_construction} & \code{.map(node2idx)} + \code{dropna} & Paires positives ou négatives dont un des deux produits n'est pas un nœud du graphe (voir H21). \\
\addlinespace
E18 & \url{prod_products_by_trigger} & \code{filter(rank <= top\_k)}, redondant avec le \code{slice(...,1,12)} déjà appliqué en amont tant qu'il n'y a qu'un seul placement par trigger. & Rien aujourd'hui --- garde-fou sans effet en placement unique ; deviendrait actif si un trigger cumulait plusieurs placements. \\

\end{longtable}
}
\refstepcounter{table}\label{tab:exclusions}
\begin{center}\footnotesize\itshape Tableau \thetable{} --- Ce qui peut être exclu du graphe.\end{center}

\end{document}
